\begin{abstract}
Artificial intelligence models are seeing widespread adoption across critical domains such as healthcare, finance, and criminal justice. Understanding the reasons behind their predictions is essential for establishing trust and reliability. This seminar paper investigates two foundational post-hoc Explainable Artificial Intelligence (XAI) methods: \textbf{Local Interpretable Model-agnostic Explanations (LIME)} and \textbf{SHapley Additive exPlanations (SHAP)}. Centered on the guiding research question: \textit{``Does a highly accurate prediction necessarily come with a reliable explanation?''}, this work covers:
(i)~\textbf{Conceptual Foundations:} Positioning post-hoc interpretability within a four-dimensional XAI taxonomy, distinguishing intrinsic from post-hoc mechanisms and local from global scopes;
(ii)~\textbf{Theoretical Formulations \& Worked Examples:} Detailing the exact optimization objectives and underlying game-theoretic axioms of LIME and SHAP, supported by step-by-step numerical calculations on an illustrative model;
(iii)~\textbf{Empirical Demonstration:} Evaluating both methods on the Breast Cancer Wisconsin dataset using a Random Forest classifier across confident, borderline, and conflicting prediction cases;
(iv)~\textbf{Runtime \& Instability:} Benchmarking computational latency across TreeSHAP, KernelSHAP, and LIME, while demonstrating explanation instability in LIME caused by stochastic sampling; and
(v)~\textbf{Faithfulness \& Limitations:} Evaluating explanation faithfulness via empirical feature-deletion experiments and addressing why post-hoc attributions must not be interpreted as causal effects.

\vspace{0.8em}
\noindent\textbf{Keywords:} Explainable AI (XAI), LIME, SHAP, KernelSHAP, TreeSHAP, Shapley Values, Local Accuracy, Consistency, Model Faithfulness.
\end{abstract}
